\chapter{Entretiens blancs}\label{ch:iv2:mock-interviews}

\begin{ivtrains}
Des entretiens complets, rédigés dans l'ordre, chacun pour un seul rôle. Chaque question donne le cadre, les règles et une feuille de score, et sa réponse contient le script. Il faut les faire d'un seul tenant, avec une horloge, sur papier et à voix haute, ou demander à un ami de lire le script en jouant le rôle du recruteur. Couvrir le script avec une feuille de papier, et la faire descendre d'une ligne d'arrêt à la fois.
\end{ivtrains}

\begin{question}[$\star\star\star\star$ \iqroles{QT} \iqfirm{teneur de marché} \iqlevel{junior}]\label{iq:iv2:mock-interviews:1}
\textbf{Trader quantitatif junior chez un teneur de marché.} Un entretien de 45 minutes avec un seul recruteur, en huit tours, sans calculatrice. Respecter le temps imparti à chaque tour, et passer au suivant dès qu'il se termine, même si l'on n'a pas fini. Dire ce que l'on fait au fur et à mesure. Chaque fois qu'on donne une estimation, donner d'abord un nombre, puis l'affiner. Après le dernier tour, noter chaque tour sur 2 d'après les notes du script, puis additionner les points.

\medskip
\begin{tabular}{@{}l|l|l|l|l@{}}
Tour & Sujet & Minutes & Réponse & Note \\ \hline
1 & Présentation & 3 & & \\[6pt]\hline
2 & Calcul & 3 & & \\[6pt]\hline
3 & Probabilités & 7 & & \\[6pt]\hline
4 & Faire un marché & 10 & & \\[6pt]\hline
5 & Dimensionner un pari & 7 & & \\[6pt]\hline
6 & Une option & 7 & & \\[6pt]\hline
7 & Jugement de trader & 5 & & \\[6pt]\hline
8 & Pression et questions au recruteur & 3 & & \\[6pt]\hline
\end{tabular}
\end{question}
\vfill\clearpage
\begin{question}[$\star\star\star\star$ \iqroles{QR} \iqfirm{hedge fund} \iqlevel{junior, confirmé}]\label{iq:iv2:mock-interviews:2}
\textbf{Chercheur quantitatif dans un hedge fund.} Un entretien de 45 minutes avec un seul chercheur, en huit tours, sur papier et sans calculatrice. Respecter le temps de chaque tour, et passer au suivant dès qu'il se termine. Dire ce que l'on fait au fur et à mesure. Après le dernier tour, noter chaque tour sur 2 d'après le script.

\medskip
\begin{tabular}{@{}l|l|l|l|l@{}}
Tour & Sujet & Minutes & Réponse & Note \\ \hline
1 & Présentation des travaux & 3 & & \\[6pt]\hline
2 & Probabilités & 7 & & \\[6pt]\hline
3 & Régression & 8 & & \\[6pt]\hline
4 & Dimensionner un signal & 7 & & \\[6pt]\hline
5 & Un backtest & 8 & & \\[6pt]\hline
6 & Coûts & 5 & & \\[6pt]\hline
7 & Données & 4 & & \\[6pt]\hline
8 & Pression et questions au recruteur & 3 & & \\[6pt]\hline
\end{tabular}
\end{question}

\vfill\clearpage
\begin{question}[$\star\star\star\star$ \iqroles{Trader} \iqfirm{banque} \iqlevel{confirmé}]\label{iq:iv2:mock-interviews:3}
\textbf{Trader sur options dans une banque.} Un entretien de 45 minutes avec le responsable d'un desk, en huit tours, sans calculatrice. Tous les taux et dividendes sont nuls. Respecter le temps de chaque tour, et passer au suivant dès qu'il se termine. Dire ce que l'on fait au fur et à mesure. Après le dernier tour, noter chaque tour sur 2 d'après le script.

\medskip
\begin{tabular}{@{}l|l|l|l|l@{}}
Tour & Sujet & Minutes & Réponse & Note \\ \hline
1 & Présentation & 2 & & \\[6pt]\hline
2 & Calcul mental & 4 & & \\[6pt]\hline
3 & Lecture d'un rapport de risque & 9 & & \\[6pt]\hline
4 & La surface de volatilité & 6 & & \\[6pt]\hline
5 & Un ordre client & 9 & & \\[6pt]\hline
6 & Un stress & 7 & & \\[6pt]\hline
7 & Jugement & 5 & & \\[6pt]\hline
8 & Questions au recruteur & 3 & & \\[6pt]\hline
\end{tabular}
\end{question}

\vfill\clearpage
\begin{question}[$\star\star\star\star$ \iqroles{Analyste} \iqfirm{hedge fund} \iqlevel{junior, confirmé}]\label{iq:iv2:mock-interviews:4}
\textbf{Analyste dans un hedge fund actions long/short.} Un entretien de 45 minutes avec un gérant de portefeuille, en huit tours, sur papier et sans calculatrice. Tous les chiffres concernent une société fictive et sont exprimés en millions de dollars, sauf par action. Respecter le temps de chaque tour, et passer au suivant dès qu'il se termine. Dire ce que l'on fait au fur et à mesure. Après le dernier tour, noter chaque tour sur 2 d'après le script.

\medskip
\begin{tabular}{@{}l|l|l|l|l@{}}
Tour & Sujet & Minutes & Réponse & Note \\ \hline
1 & Présentation & 2 & & \\[6pt]\hline
2 & Du chiffre d'affaires à la trésorerie & 7 & & \\[6pt]\hline
3 & Une valorisation & 8 & & \\[6pt]\hline
4 & Sensibilité & 6 & & \\[6pt]\hline
5 & Une opération & 6 & & \\[6pt]\hline
6 & Une idée de vente à découvert & 8 & & \\[6pt]\hline
7 & Questions insistantes & 5 & & \\[6pt]\hline
8 & Questions au recruteur & 3 & & \\[6pt]\hline
\end{tabular}
\end{question}

\vfill\clearpage
\begin{question}[$\star\star\star\star$ \iqroles{DS} \iqfirm{crypto} \iqlevel{junior, confirmé}]\label{iq:iv2:mock-interviews:5}
\textbf{Data scientist dans une société crypto.} Un entretien de 45 minutes avec un data scientist senior, en huit tours, sur papier et sans calculatrice. Respecter le temps de chaque tour, et passer au suivant dès qu'il se termine. Dire ce que l'on fait au fur et à mesure. Après le dernier tour, noter chaque tour sur 2 d'après le script.

\medskip
\begin{tabular}{@{}l|l|l|l|l@{}}
Tour & Sujet & Minutes & Réponse & Note \\ \hline
1 & Présentation des travaux & 3 & & \\[6pt]\hline
2 & Une expérience & 7 & & \\[6pt]\hline
3 & Un événement rare & 7 & & \\[6pt]\hline
4 & Fuite d'information & 6 & & \\[6pt]\hline
5 & Données corrompues & 6 & & \\[6pt]\hline
6 & Nombreuses variables & 7 & & \\[6pt]\hline
7 & Rendements à court horizon & 6 & & \\[6pt]\hline
8 & Questions au recruteur & 3 & & \\[6pt]\hline
\end{tabular}
\end{question}

\vfill\clearpage
\begin{question}[$\star\star\star\star$ \iqroles{QD} \iqfirm{banque} \iqlevel{junior, confirmé}]\label{iq:iv2:mock-interviews:6}
\textbf{Développeur quantitatif dans une banque.} Un entretien de 45 minutes avec un développeur senior, en huit tours, sur papier et sans calculatrice. On décrit le code en mots~: personne ne demande de l'écrire ici. Respecter le temps de chaque tour, et passer au suivant dès qu'il se termine. Dire ce que l'on fait au fur et à mesure. Après le dernier tour, noter chaque tour sur 2 d'après le script.

\medskip
\begin{tabular}{@{}l|l|l|l|l@{}}
Tour & Sujet & Minutes & Réponse & Note \\ \hline
1 & Présentation & 2 & & \\[6pt]\hline
2 & Erreur de Monte-Carlo & 7 & & \\[6pt]\hline
3 & Grilles et convergence & 7 & & \\[6pt]\hline
4 & Virgule flottante & 6 & & \\[6pt]\hline
5 & Risque à grande échelle & 8 & & \\[6pt]\hline
6 & Perturbations et bogues & 7 & & \\[6pt]\hline
7 & Une panne en production & 5 & & \\[6pt]\hline
8 & Questions au recruteur & 3 & & \\[6pt]\hline
\end{tabular}
\end{question}

\vfill\clearpage
\begin{question}[$\star\star\star\star$ \iqroles{Dev} \iqfirm{teneur de marché} \iqlevel{junior, confirmé}]\label{iq:iv2:mock-interviews:7}
\textbf{Développeur logiciel chez un teneur de marché.} Un entretien de 45 minutes avec un ingénieur confirmé, en huit tours, sur papier et sans calculatrice. Vous décrivez le code à l'oral~; personne ne vous demande de l'écrire ici. Tenez-vous-en au temps imparti pour chaque tour, et passez au suivant quand il se termine. Dites ce que vous faites au fur et à mesure. À l'issue du dernier tour, notez chaque tour sur 2 d'après le script.

\medskip
\begin{tabular}{@{}l|l|l|l|l@{}}
Tour & Sujet & Minutes & Réponse & Note \\ \hline
1 & Présentation & 2 & & \\[6pt]\hline
2 & Une structure de données & 8 & & \\[6pt]\hline
3 & Hachage & 6 & & \\[6pt]\hline
4 & Une file d'attente & 7 & & \\[6pt]\hline
5 & Quel volume de données & 5 & & \\[6pt]\hline
6 & Concurrence & 7 & & \\[6pt]\hline
7 & Un pic & 7 & & \\[6pt]\hline
8 & Questions au recruteur & 3 & & \\[6pt]\hline
\end{tabular}
\end{question}

\vfill\clearpage
\begin{question}[$\star\star\star\star$ \iqroles{LLE} \iqfirm{teneur de marché} \iqlevel{confirmé}]\label{iq:iv2:mock-interviews:8}
\textbf{Ingénieur basse latence chez un teneur de marché.} Un entretien de 45 minutes avec un ingénieur basse latence, en huit tours, sur papier et sans calculatrice. Les ordres de grandeur du matériel sont donnés dans les prompts, à titre d'hypothèses pour l'entretien. Tenez-vous-en au temps imparti pour chaque tour, et passez au suivant quand il se termine. Dites ce que vous faites au fur et à mesure. À l'issue du dernier tour, notez chaque tour sur 2 d'après le script.

\medskip
\begin{tabular}{@{}l|l|l|l|l@{}}
Tour & Sujet & Minutes & Réponse & Note \\ \hline
1 & Présentation & 2 & & \\[6pt]\hline
2 & Budget en cycles & 5 & & \\[6pt]\hline
3 & Câblage et distance & 7 & & \\[6pt]\hline
4 & Une chaîne d'étapes & 6 & & \\[6pt]\hline
5 & Branchements et caches & 7 & & \\[6pt]\hline
6 & Le noyau & 6 & & \\[6pt]\hline
7 & Mesurer un changement & 9 & & \\[6pt]\hline
8 & Questions au recruteur & 3 & & \\[6pt]\hline
\end{tabular}
\end{question}

\vfill\clearpage
\begin{question}[$\star\star\star\star$ \iqroles{Risque} \iqfirm{banque} \iqlevel{junior, confirmé}]\label{iq:iv2:mock-interviews:9}
\textbf{Responsable du risque de marché dans une banque.} Un entretien de 45 minutes avec un directeur des risques, en huit tours, sur papier et sans calculatrice. Traitez les rendements comme gaussiens, sauf si un prompt indique le contraire. Tenez-vous-en au temps imparti pour chaque tour, et passez au suivant quand il se termine. Dites ce que vous faites au fur et à mesure. À l'issue du dernier tour, notez chaque tour sur 2 d'après le script.

\medskip
\begin{tabular}{@{}l|l|l|l|l@{}}
Tour & Sujet & Minutes & Réponse & Note \\ \hline
1 & Présentation & 2 & & \\[6pt]\hline
2 & VaR & 7 & & \\[6pt]\hline
3 & Deux desks & 6 & & \\[6pt]\hline
4 & Un backtest de VaR & 8 & & \\[6pt]\hline
5 & Mesures de queue de distribution & 5 & & \\[6pt]\hline
6 & Liquidité & 6 & & \\[6pt]\hline
7 & Crédit & 4 & & \\[6pt]\hline
8 & Un trader en dépassement, et vos questions & 7 & & \\[6pt]\hline
\end{tabular}
\end{question}

\vfill\clearpage
\begin{question}[$\star\star\star\star$ \iqroles{QT, Trader} \iqfirm{crypto} \iqlevel{confirmé}]\label{iq:iv2:mock-interviews:10}
\textbf{Dernier tour~: trader confirmé dans une société crypto.} Un dernier tour de 60 minutes avec le responsable du trading, en huit tours, sur papier et sans calculatrice. Tout ce qui figure ici est inventé, aucune plateforme n'est réelle, et rien sur les règles d'une plateforme n'est présenté comme un fait. Tenez-vous-en au temps imparti pour chaque tour, et passez au suivant quand il se termine. Dites ce que vous faites au fur et à mesure. À l'issue du dernier tour, notez chaque tour sur 2 d'après le script.

\medskip
\begin{tabular}{@{}l|l|l|l|l@{}}
Tour & Sujet & Minutes & Réponse & Note \\ \hline
1 & Votre parcours & 5 & & \\[6pt]\hline
2 & Un basis trade & 10 & & \\[6pt]\hline
3 & Liquidation & 8 & & \\[6pt]\hline
4 & Risque de plateforme & 7 & & \\[6pt]\hline
5 & Deux plateformes & 8 & & \\[6pt]\hline
6 & Dimensionnement & 8 & & \\[6pt]\hline
7 & Une mauvaise semaine & 9 & & \\[6pt]\hline
8 & Questions au recruteur & 5 & & \\[6pt]\hline
\end{tabular}
\end{question}

